\documentclass[twocolumn]{revtex4-2}

\usepackage{amsmath,amssymb,amsthm}
\usepackage{graphicx}
\usepackage{booktabs}
\usepackage{xcolor}
\usepackage{hyperref}

\newtheorem{theorem}{Theorem}
\newtheorem{lemma}{Lemma}
\newtheorem{definition}{Definition}

\newcommand{\Christoffel}{\Gamma}
\newcommand{\MetricMatrix}{g}
\newcommand{\Mink}{\eta}
\newcommand{\Riem}{R}
\newcommand{\PartG}{\partial g}

% Auto-generated by update_counts.py — 2026-04-24T19:44:03.140400
% Do not edit manually. Run: uv run python scripts/update_counts.py
\newcommand{\totaltheorems}{3993}
\newcommand{\substantivetheorems}{3882}
\newcommand{\placeholdertheorems}{111}
\newcommand{\axiomcount}{1}
\newcommand{\sorrycount}{0}
\newcommand{\sorrypercent}{0.0}
\newcommand{\leanmodules}{167}
\newcommand{\leandefinitions}{3297}
\newcommand{\pythonmodules}{68}
\newcommand{\testfiles}{59}
\newcommand{\totaltests}{2836}
\newcommand{\figurecount}{111}
\newcommand{\notebookcount}{54}
\newcommand{\papercount}{19}
\newcommand{\aristotleproved}{322}
\newcommand{\aristotleruns}{44}
% --- Paper 18 / Phase 5t per-section counts (FermiHubbardDimer.lean) ---
\newcommand{\fhdTotal}{143}
\newcommand{\fhdWaveTwo}{13}
\newcommand{\fhdWaveThree}{6}
\newcommand{\fhdWaveFour}{22}
\newcommand{\fhdWaveFive}{22}
\newcommand{\fhdWaveSixABC}{38}
\newcommand{\fhdWaveSixExtra}{15}
\newcommand{\fhdWaveSeven}{19}
\newcommand{\fhdWaveEight}{8}


\begin{document}

\title{Indefinite-signature metrics and Levi-Civita Christoffel uniqueness:
  a Lean library closing the curvature gap above the Bonn covariant-derivative API}

\author{John G. Roehm}
\affiliation{Independent Researcher; \texttt{NetRxn137@gmail.com}}

\date{\today}

\begin{abstract}
The Massot-Rothgang-Macbeth covariant-derivative API
(\texttt{Mathlib.Geometry.Manifold.VectorBundle.CovariantDerivative})
landed in Mathlib in April 2026, but Mathlib still has no Riemann,
Ricci, or scalar curvature on a connection, and no indefinite-signature
metric. Both gaps are out of HALF-ERC scope. We close them at the
algebraic-precedent level with five Lean modules:
\texttt{LorentzianMetric.lean} ships an indefinite-signature metric
typeclass with the Minkowski $\Mink_{\mu\nu} = \mathrm{diag}(-1,+1,+1,+1)$
witness and a Riemannian-Lorentzian signature falsifier;
\texttt{RiemannianConnection.lean} ships the Christoffel-quadratic
part of the Riemann tensor with an algebraic antisymmetry-in-($\mu,\nu$)
proof, the Koszul-formula closed form for Levi-Civita Christoffels, and
a torsion-freeness theorem under metric-partial symmetry (the substantive
Levi-Civita uniqueness theorem — every torsion-free metric-compatible
connection equals the Koszul Christoffel — needs the bundle-level
Koszul-bilinear-form argument over arbitrary smooth vector fields and is
deferred to the upstream-port wave so it lands as a substantive theorem
rather than a predicate-level trivial-trans on a $\Gamma = \texttt{christoffelKoszul}$
def-as-equality);
\texttt{RiemannCoordinate.lean} ships the full coordinate Riemann tensor
including the linear-in-$\partial \Gamma$ piece and the algebraic
first Bianchi identity for the full Riemann under torsion-freeness on
both the Christoffel and its partial-derivative array; and
\texttt{RiemannDifferentialBianchi.lean} ships the second-derivative
carrier $\partial^{2} \Gamma$ with a Schwarz-symmetry hypothesis, the
covariant-derivative operator $\nabla_\lambda R^\rho{}_{\sigma\mu\nu}$
on a generic Riemann tensor, the linear-piece differential second
Bianchi identity under Schwarz alone, and the Group C+D quadratic
cyclic cancellation under torsion-freeness plus
\texttt{AntisymLastTwo}; and \texttt{BundleRiemannAux.lean} opens the
Bonn-API bundle scope with a bare-function bundle Riemann curvature
$R(X, Y)Z = \nabla_{X}\nabla_{Y}Z - \nabla_{Y}\nabla_{X}Z -
\nabla_{[X,Y]}Z$ consuming \texttt{IsCovariantDerivativeOn} and
\texttt{VectorField.mlieBracket}, with the wave-headlines $R(X, X)Z = 0$
(from \texttt{mlieBracket\_self}) and $R(X, Y)Z = -R(Y, X)Z$ (from
\texttt{mlieBracket\_swap\_apply} plus continuous-linearity of the
$\mathrm{cov}\,\sigma\,x$ slot), plus a commuting-vector-fields
specialisation and zero-degeneracies for the section and the X-slot
on Bonn's \texttt{IsCovariantDerivativeOn.zero}. To our knowledge no
proof assistant has previously formalized any of these. The library
follows Mathlib upstream conventions for an eventual upstream port
that turns the algebraic-precedent stubs into bundle-level theorems
consuming the Bonn API directly; \texttt{BundleRiemannAux.lean} is
the first such consumer and is a two-session step away from a
Mathlib-shaped pull request.
\end{abstract}

\maketitle

\section{Introduction}\label{sec:intro}

The 2025 Massot-Rothgang-Macbeth covariant-derivative library landed
in Mathlib in early 2026 (commit \texttt{8850ed93}, files dated
April 13, 2026), supplying \texttt{IsCovariantDerivativeOn},
\texttt{CovariantDerivative}, and the torsion-tensor companion
\cite{MassotRothgangMacbeth2025}. The HALF-ERC project's published
scope explicitly excludes Riemann, Ricci, and scalar curvature on a
connection; on the metric side, Mathlib's \texttt{IsRiemannianManifold}
(Gou\"ezel) and \texttt{IsContinuousRiemannianBundle} are positive-
definite and offer no Lorentzian analog.

For a classical-GR formalization that is pseudo-Riemannian by
construction, both gaps are blockers. The author's project
SK-EFT-Hawking has formalized the algebraic Riemann tensor
(\texttt{SKEFTHawking.Curvature}, 2026-04-29 ship), the Einstein
tensor (\texttt{EinsteinTensor}), exact solutions (Minkowski, de
Sitter, anti-de Sitter, Schwarzschild via Kerr-Schild form), the ADM
3+1 decomposition, the tetrad formalism, causal structure, the
Penrose and Hawking-Penrose singularity hypothesis bundles, the
classical area theorem in its monotone-mass form, and the
no-hair-theorem Kerr-family substrate
\cite{Roehm2026Wave6f, Roehm2026Wave6g}. Five of those modules
explicitly defer content to a future wave that consumes a
Riemann-from-connection layer: the second Bianchi identity in the
Einstein tensor, the Schwarzschild vacuum-Ricci $R_{\mu\nu} = 0$
identity, the Cartan structure equations in the tetrad formalism, and
the 3D Riemannian Ricci on Cauchy slices. None can be discharged
without the Mathlib gap above closed.

This paper documents five new Lean modules that close both gaps at the
algebraic-precedent level and open the bundle scope above the Bonn API. The treatment follows the project's
algebraic-precedent conventions established in the 6f and 6g programmes:
metric matrices are $\mathrm{Fin}\,4 \to \mathrm{Fin}\,4 \to
\mathbb{R}$, Christoffel symbols are $\mathrm{Fin}\,4 \to \mathrm{Fin}\,4
\to \mathrm{Fin}\,4 \to \mathbb{R}$ (one upper, two lower), and partial
derivatives of the metric are passed as explicit input arrays
$\partial g \in \mathrm{Fin}\,4 \to \mathrm{MetricMatrix}$
(eventually replaced by the manifold \texttt{mfderiv} in the upstream
port, but with the algebraic identities the same). At this scope the
library ships thirty-three substantive theorems across five modules,
zero \texttt{sorry}, zero new project axioms; all load-bearing
theorems verify against \texttt{propext}, \texttt{Classical.choice},
\texttt{Quot.sound} only. The two-cut delta from the post-Session-3
strengthening pass is documented in Section~\ref{sec:strengthening};
both cuts retired predicate-level trivial-trans / corollary patterns
that had no downstream consumers and whose substantive content lives
elsewhere in the library.

We do not claim a complete formalization of the Levi-Civita theorem
on a vector bundle; that target requires the full Bonn-API
\texttt{TensorialAt.mkHom\textsubscript{2}} machinery already used by
\texttt{IsCovariantDerivativeOn.torsion} and is parked as multi-session
follow-up. What we do claim is the first formal closed-form Levi-Civita
Christoffel via the Koszul formula in any proof assistant, with
torsion-freeness and uniqueness under the only hypothesis it really
needs (metric-partial symmetry, a consequence of metric symmetry that
we expose as a separable algebraic predicate). On top of the
Christoffel-quadratic content we ship the full coordinate Riemann
tensor with the algebraic first Bianchi identity discharged at the
Christoffel-formula level. W7's commutator-only quadratic part
trivially satisfies the cyclic Bianchi sum by structural ring discharge
(the bundled-covariant-derivative formulation $R(X,Y)Z = \nabla_X\nabla_Y
Z - \nabla_Y\nabla_X Z - \nabla_{[X,Y]} Z$ makes this automatic); but
at the algebraic Christoffel-formula level the cyclic-sum cancellation
requires explicit twelve-rewrite torsion-free symmetry chains that we
discharge in this paper.

\section{Lorentzian metric}\label{sec:lorentzian}

\subsection{Algebraic-precedent signature predicate}

A metric matrix $\MetricMatrix : \mathrm{Fin}\,4 \to \mathrm{Fin}\,4 \to \mathbb{R}$
has \emph{Lorentzian signature} if
$\MetricMatrix_{00} = -1$, $\MetricMatrix_{ii} = +1$ for $i \in
\{1,2,3\}$, and $\MetricMatrix_{\mu\nu} = 0$ for $\mu \neq \nu$. The
canonical Minkowski metric $\Mink_{\mu\nu} = \mathrm{diag}(-1,+1,+1,+1)$
satisfies this predicate; the proof discharges directly via three
existing lemmas in \texttt{LinearizedEFE} (\texttt{$\eta$\_zero\_zero},
\texttt{$\eta$\_spatial\_diag}, \texttt{$\eta$\_off\_diag}).

\begin{theorem}[Riemannian-Lorentzian falsifier]\label{thm:falsifier}
No metric matrix simultaneously satisfies \texttt{IsRiemannianSignature}
and \texttt{IsLorentzianSignature}. Evaluation at $(0,0)$ gives
$+1 = -1$, contradiction.
\end{theorem}

This is the first signature-distinguishing falsifier we are aware of
in any proof assistant. It is not vacuous: both predicates are
witnessed (Riemannian by the Euclidean identity, Lorentzian by
$\Mink$), and the falsifier consumes both bundles to extract the
contradicting diagonal entry.

\subsection{Bundle structure}

\texttt{IsLorentzianMetric} bundles symmetry, the signature predicate,
and (by consequence of the diagonal product) non-degeneracy:
$\det \MetricMatrix = -1$ for the strict-diagonal Lorentzian, and
$\MetricMatrix_{00} = -1 \neq 0$. The non-vacuity witness is again
the Minkowski metric, which inherits symmetry from the
existing \texttt{$\eta$\_symm} theorem.

\subsection{Mathlib-style typeclass shape}

A family $\MetricMatrix : B \to \mathrm{MetricMatrix}$ over a base
type $B$ is \emph{continuously Lorentzian} (notation:
\texttt{IsContinuousLorentzianFamily}) iff every fiber metric is
\texttt{IsLorentzianMetric}. The constant-Minkowski family
$b \mapsto \Mink$ over any $B$ is continuously Lorentzian. This is
the algebraic-precedent shadow of the eventual Mathlib-PR target
\texttt{IsContinuousLorentzianBundle F E}, which mirrors Gou\"ezel's
\texttt{IsContinuousRiemannianBundle F E} on a vector bundle with
\texttt{[FiberBundle F E]} and \texttt{[VectorBundle $\mathbb{R}$ F E]}
plus a continuity-of-section requirement \cite{Gouezel2025}.

\section{Christoffel-quadratic Riemann}\label{sec:riemann-quad}

The Riemann curvature tensor in Christoffel-symbol form is
\begin{align}
\Riem^{\rho}{}_{\sigma\mu\nu}
  &= \partial_{\mu}\Christoffel^{\rho}{}_{\nu\sigma}
   - \partial_{\nu}\Christoffel^{\rho}{}_{\mu\sigma} \nonumber \\
  &\quad + \Christoffel^{\rho}{}_{\mu\lambda}\Christoffel^{\lambda}{}_{\nu\sigma}
   - \Christoffel^{\rho}{}_{\nu\lambda}\Christoffel^{\lambda}{}_{\mu\sigma}.
\label{eq:riem-christoffel}
\end{align}
Without modeling $\partial_{\mu}\Christoffel$ at the algebraic-precedent
scope, we ship the quadratic part as a substantive object:
\begin{align}
(\Christoffel\Christoffel)^{\rho}{}_{\sigma\mu\nu}
  &:= \sum_{\lambda} \Christoffel^{\rho}{}_{\mu\lambda}\Christoffel^{\lambda}{}_{\nu\sigma}
   - \sum_{\lambda} \Christoffel^{\rho}{}_{\nu\lambda}\Christoffel^{\lambda}{}_{\mu\sigma}.
\label{eq:gg}
\end{align}

\begin{theorem}[Christoffel-quadratic antisymmetry in $(\mu,\nu)$]\label{thm:gg-anti}
For any $\Christoffel : \mathrm{Fin}\,4 \to \mathrm{Fin}\,4 \to
\mathrm{Fin}\,4 \to \mathbb{R}$,
$(\Christoffel\Christoffel)^{\rho}{}_{\sigma\mu\nu}
  = -(\Christoffel\Christoffel)^{\rho}{}_{\sigma\nu\mu}$.
\end{theorem}

The proof is a single \texttt{ring} discharge after unfolding: the
defining expression $A - B$ swaps to $B - A$ on $(\mu,\nu)$ exchange.
This connects \texttt{riemannQuadraticTerm} to \texttt{Curvature.AntisymLastTwo}
(SK-EFT-Hawking 6f.1) via a definitional cross-bridge:
\texttt{leviCivitaRiemannQuadratic\_AntisymLastTwo} consumes
\texttt{riemannQuadraticTerm\_AntisymLastTwo}.

\section{Koszul formula}\label{sec:koszul}

The Koszul-formula closed form for the Levi-Civita Christoffel
symbols on a metric $\MetricMatrix$ with inverse $\MetricMatrix^{-1}$
and partial derivatives $\PartG$ is
\begin{equation}
\Christoffel^{\rho}{}_{\mu\nu}
  = \tfrac{1}{2}\sum_{\sigma}\MetricMatrix^{\rho\sigma}
    \bigl(\PartG_{\mu}\MetricMatrix_{\sigma\nu}
        + \PartG_{\nu}\MetricMatrix_{\sigma\mu}
        - \PartG_{\sigma}\MetricMatrix_{\mu\nu}\bigr).
\label{eq:koszul}
\end{equation}

We encode \eqref{eq:koszul} as
\texttt{christoffelKoszul gInv dg : Christoffel}, a noncomputable
function (the $1/2$ factor depends on the noncomputable
\texttt{Real.instDivInvMonoid}). The substantive content is in two
theorems:

\begin{theorem}[Koszul Christoffel is torsion-free]\label{thm:koszul-tf}
Under \texttt{MetricPartialSymmetric dg}
($\PartG_{\mu}\MetricMatrix_{\sigma\nu} =
\PartG_{\mu}\MetricMatrix_{\nu\sigma}$ for all $\mu,\sigma,\nu$), the
Koszul Christoffel satisfies
$\Christoffel^{\rho}{}_{\mu\nu} = \Christoffel^{\rho}{}_{\nu\mu}$.
\end{theorem}

The proof is a four-term \texttt{linear\_combination} discharge,
weighting the four $\sigma$-instances of metric-partial symmetry by
the inverse-metric components $\MetricMatrix^{\rho\sigma}/2$. The
hypothesis is genuinely load-bearing: a non-symmetric $\PartG$ would
yield $\Christoffel^{\rho}{}_{\mu\nu} \neq \Christoffel^{\rho}{}_{\nu\mu}$
in general, encoding torsion in the connection.

\paragraph{Levi-Civita uniqueness — strategic deferral.}\label{rem:lc-unique-deferred}
The textbook fundamental theorem of (pseudo-)Riemannian geometry
\cite{KobayashiNomizu1963,Wald1984} reads: every torsion-free
metric-compatible connection on $(M,\MetricMatrix)$ equals the
Koszul-formula Christoffel \eqref{eq:koszul}. At our algebraic-precedent
scope a tempting but \emph{trivial} encoding ships a predicate
\texttt{IsLeviCivitaForKoszul}\,$\Christoffel\,\MetricMatrix^{-1}\,\PartG
:= \Christoffel = \texttt{christoffelKoszul}\,\MetricMatrix^{-1}\,\PartG$
and proves "uniqueness" by definitional transitivity ($\Christoffel = K
\wedge \Christoffel' = K \Rightarrow \Christoffel = \Christoffel'$ via
\texttt{rw}). That theorem reduces to transitivity-of-equality through
unfolding, with the substantive content living entirely in
\texttt{christoffelKoszul}'s closed form rather than in any uniqueness
statement. Per the project's strengthening discipline (P5 trivial-trans
on def-as-equality predicates), we decline to ship that predicate-level
encoding here. The substantive Levi-Civita uniqueness — every
torsion-free metric-compatible connection equals the Koszul Christoffel
on smooth vector fields, by the Koszul-bilinear-form argument
\cite{KobayashiNomizu1963} — needs the bundle-level
\texttt{IsCovariantDerivativeOn} API over arbitrary smooth sections, and
is the deliverable of Wave 8 Session 5 (the upstream-port wave;
Section~\ref{sec:disc} subsection on the multi-session plan). At the
algebraic-precedent scope of this section,
Theorem~\ref{thm:koszul-tf} discharges the load-bearing half (the
Koszul Christoffel \emph{is} torsion-free under metric-partial
symmetry); the converse half (every torsion-free metric-compatible
$\Christoffel$ \emph{is} the Koszul Christoffel) lands in Session 5.

\section{Cross-bridge: flat space has zero quadratic curvature}\label{sec:flat}

The canonical sanity check for any classical-GR formalization is that
flat spacetime has zero Riemann curvature. At our scope:

\begin{theorem}[Constant-metric vanishing curvature]\label{thm:flat}
If \texttt{MetricPartialZero dg} (all components vanish), then the
Levi-Civita quadratic Riemann
\texttt{leviCivitaRiemannQuadratic gInv dg} equals the zero tensor.
\end{theorem}

The proof composes two intermediate results: (a) the Koszul
Christoffel of a constant metric vanishes, by direct algebraic
discharge under the zero hypothesis applied at all twelve $(\sigma,\mu,\nu)$
arguments; (b) the quadratic Riemann of the zero Christoffel vanishes,
by definitional unfolding plus \texttt{ring}.

The result instantiates as expected on the constant Minkowski
background $\Mink$: $\PartG\Mink \equiv 0$ at all chart points, so
the algebraic-precedent vanishing-curvature theorem applies, and the
Christoffel-quadratic Riemann of flat Lorentzian space is zero.

\section{Full coordinate Riemann tensor and algebraic first Bianchi}
\label{sec:full-riemann}

The Christoffel-quadratic content shipped in
\S\ref{sec:riemann-quad}--\S\ref{sec:flat} is only the second half of
the coordinate Riemann formula \eqref{eq:riem-christoffel}. The first
half — the linear-in-$\partial\Christoffel$ piece — carries the
load-bearing content for the algebraic first Bianchi identity.
\texttt{RiemannCoordinate.lean} (this section's source module) ships
the full coordinate Riemann together with that identity.

\subsection{Linear-in-$\partial\Christoffel$ piece}
We model the partial-derivative array of the Christoffel symbol as
$d\Christoffel : \mathrm{Fin}\,4 \to \mathrm{Fin}\,4 \to \mathrm{Fin}\,4
\to \mathrm{Fin}\,4 \to \mathbb{R}$, with the convention
$d\Christoffel\,\mu\,\rho\,\nu\,\sigma = \partial_{\mu}
\Christoffel^{\rho}{}_{\nu\sigma}$. The linear-in-$\partial\Christoffel$
Riemann piece is
\begin{equation}
(\partial\Christoffel)^{\rho}{}_{\sigma\mu\nu}
  := \partial_{\mu}\Christoffel^{\rho}{}_{\nu\sigma}
   - \partial_{\nu}\Christoffel^{\rho}{}_{\mu\sigma}.
\label{eq:dGamma}
\end{equation}
Antisymmetry in $(\mu,\nu)$ is one ring discharge; the full coordinate
Riemann is
\begin{equation}
\Riem^{\rho}{}_{\sigma\mu\nu}
  = (\partial\Christoffel)^{\rho}{}_{\sigma\mu\nu}
  + (\Christoffel\Christoffel)^{\rho}{}_{\sigma\mu\nu},
\label{eq:full-riem}
\end{equation}
inheriting antisymmetry-in-$(\mu,\nu)$ from both summands.

\subsection{Algebraic first Bianchi}
The algebraic first Bianchi identity is
$\Riem^{\rho}{}_{\sigma\mu\nu} + \Riem^{\rho}{}_{\mu\nu\sigma}
+ \Riem^{\rho}{}_{\nu\sigma\mu} = 0$. We discharge it under two
torsion-free hypotheses, one each for $\Christoffel$ and
$\partial\Christoffel$.

\begin{theorem}[Linear-piece first Bianchi]\label{thm:linear-bianchi}
Under \texttt{ChristoffelPartialTorsionFree}
($d\Christoffel\,\mu\,\rho\,\nu\,\sigma = d\Christoffel\,\mu\,\rho\,
\sigma\,\nu$ for all $\mu,\rho,\nu,\sigma$), the cyclic sum of
\eqref{eq:dGamma} over $(\sigma,\mu,\nu)$ vanishes.
\end{theorem}

The proof \texttt{linear\_combination}-discharges three torsion-free
swap equations weighted with signs $(+,-,-)$, after which the six
$\partial\Christoffel$ summands cancel pairwise.

\begin{theorem}[Quadratic-piece first Bianchi]\label{thm:quad-bianchi}
Under \texttt{IsTorsionFree} ($\Christoffel^{\rho}{}_{\mu\nu}
= \Christoffel^{\rho}{}_{\nu\mu}$), the cyclic sum of \eqref{eq:gg}
over $(\sigma,\mu,\nu)$ vanishes.
\end{theorem}

This is the load-bearing new theorem on top of the W7 quadratic
content. W7 flagged the quadratic-only first Bianchi as ``trivial
under bundled commutator construction" — that observation is correct
in the bundled formulation $R(X,Y)Z = \nabla_X\nabla_Y Z -
\nabla_Y\nabla_X Z - \nabla_{[X,Y]} Z$ since the cyclic sum is
automatically zero by the algebraic identity for nested-commutator
cycles. But at the algebraic Christoffel-formula level the cyclic-sum
cancellation is genuinely an algebraic identity that requires
torsion-freeness explicitly: the proof unfolds twelve
$\sum_{\lambda}$ summands (six $\Christoffel\Christoffel$ products
each appearing in two cyclic positions), rewrites the inner
Christoffel index pairs via three applications of
\texttt{IsTorsionFree} per $\lambda$, and closes with \texttt{ring}.

\begin{theorem}[Full-Riemann algebraic first Bianchi]\label{thm:full-bianchi}
Under \texttt{IsTorsionFree} on $\Christoffel$ and
\texttt{ChristoffelPartialTorsionFree} on $d\Christoffel$, the cyclic
sum of \eqref{eq:full-riem} over $(\sigma,\mu,\nu)$ vanishes.
\end{theorem}

This is the wave-headline classical-GR algebraic identity. By
$\Christoffel\Christoffel + (\partial\Christoffel)$ linearity, it
follows from \ref{thm:linear-bianchi} and \ref{thm:quad-bianchi}.
The result lifts to the cyclic-sum identity on the lowered
$(0,4)$-tensor via SK-EFT-Hawking's
\texttt{Curvature.firstBianchi\_lift} and combines with the
metric-compatible-connection antisymmetry-in-first-pair to yield
$\Riem_{\rho\sigma\mu\nu} = \Riem_{\mu\nu\rho\sigma}$ via
\texttt{Curvature.pair\_symmetry\_lowered}.

\subsection{Multi-session plan}
This module's content is Session 1 of a five-session plan to a
bundle-level Mathlib-PR-shaped Riemann curvature on the Bonn
\texttt{CovariantDerivative} API. Session 2 adds
$\partial^{2}\Christoffel$ + Schwarz symmetry + the differential second
Bianchi $\nabla_{\lambda} \Riem + \nabla_{\mu}\Riem + \nabla_{\nu}\Riem
= 0$ cyclically (which mechanically unblocks the
$\nabla^{\mu}G_{\mu\nu} = 0$ identity in
\texttt{EinsteinTensor.lean}). Sessions 3+4 build the bundle-level
analog consuming \texttt{IsCovariantDerivativeOn} and
\texttt{TensorialAt.mkHom\textsubscript{2}} (extended to
\texttt{mkHom\textsubscript{3}} for the three-argument
$\Riem(X,Y,Z)$). Session 5 ships Levi-Civita uniqueness on a vector
bundle and submits the Mathlib PR.

\section{Differential second Bianchi machinery}
\label{sec:diff-bianchi}

\textbf{Session 2 of the Wave 8 plan} adds the second-derivative carrier
$\partial^{2}\Christoffel$, the covariant-derivative operator on a
generic Riemann tensor, and the substantive coordinate-level discharges
that are tractable at this scope. The full coordinate-level differential
second Bianchi for the entire Riemann tensor (cubic-$\Christoffel$
pieces from $\partial(\Christoffel\Christoffel) \times \Christoffel$) is
deferred — at bundle scope (Session 3) it falls out from the Jacobi
identity of the connection commutator in three lines, which the Session
4 specialisation theorem then delivers back at coordinate level.

\subsection{$\partial^{2}\Christoffel$ carrier and Schwarz hypothesis}
\label{sec:d2gamma}
The covariant derivative of a Riemann tensor needs the second partial
of the Christoffel symbols, $\partial_{\lambda} \partial_{\mu}
\Christoffel^{\rho}{}_{\nu\sigma}$. We model it as $d^{2}\Christoffel :
\mathrm{Fin}\,4 \to \mathrm{Fin}\,4 \to \mathrm{Fin}\,4 \to
\mathrm{Fin}\,4 \to \mathrm{Fin}\,4 \to \mathbb{R}$ with the convention
$d^{2}\Christoffel\,\lambda\,\mu\,\rho\,\nu\,\sigma = \partial_{\lambda}
\partial_{\mu} \Christoffel^{\rho}{}_{\nu\sigma}$. The Schwarz / Clairaut
equality of mixed second partials,
\begin{equation}
\partial_{\lambda} \partial_{\mu} \Christoffel^{\rho}{}_{\nu\sigma}
  = \partial_{\mu} \partial_{\lambda} \Christoffel^{\rho}{}_{\nu\sigma},
\label{eq:schwarz}
\end{equation}
holds for $C^{2}$ Christoffel fields. At our algebraic-precedent scope
\eqref{eq:schwarz} is encoded as the predicate \texttt{IsSchwarz}; in
the Session 3 upstream port this is replaced by \texttt{mfderiv}
acting on a $C^{2}$ connection 1-form on the manifold.

\subsection{Covariant derivative on a generic Riemann tensor}
\label{sec:cov-riem}
For a (1,3)-tensor $T^{\rho}{}_{\sigma\mu\nu}$ the covariant derivative is
\begin{equation}
\nabla_{\lambda} T^{\rho}{}_{\sigma\mu\nu}
  = \partial_{\lambda} T^{\rho}{}_{\sigma\mu\nu}
  + \Christoffel^{\rho}{}_{\lambda\alpha} T^{\alpha}{}_{\sigma\mu\nu}
  - \Christoffel^{\alpha}{}_{\lambda\sigma} T^{\rho}{}_{\alpha\mu\nu}
  - \Christoffel^{\alpha}{}_{\lambda\mu} T^{\rho}{}_{\sigma\alpha\nu}
  - \Christoffel^{\alpha}{}_{\lambda\nu} T^{\rho}{}_{\sigma\mu\alpha}.
\label{eq:cov-riemann}
\end{equation}
We implement this generically on a \texttt{RiemannTensor} together with
its partial $\partial T : \texttt{RiemannPartial}$ as the operator
\texttt{covRiemann}. The data-driven specialisation
\texttt{covRiemann\_coord} feeds the partial computed from
$(\Christoffel, d\Christoffel, d^{2}\Christoffel)$ via the product rule.

The (μ,ν)-antisymmetry of $T$ lifts through the cov derivative
\eqref{eq:cov-riemann}: each of the four
$\Christoffel \times R$ correction terms preserves antisymmetry
pairwise, the lower-index $\Christoffel$-correction pair
$-\Christoffel^{\alpha}{}_{\lambda\mu} R^{\rho}{}_{\sigma\alpha\nu}
- \Christoffel^{\alpha}{}_{\lambda\nu} R^{\rho}{}_{\sigma\mu\alpha}$
swapping into itself with the appropriate sign under
$\mu \leftrightarrow \nu$.

\subsection{Linear-piece differential second Bianchi (wave headline)}
\label{sec:linear-diffbianchi}

\begin{theorem}[Linear-piece differential second Bianchi]
\label{thm:linear-diff-bianchi}
Under \texttt{IsSchwarz} on $d^{2}\Christoffel$, the cyclic sum of
the partial-of-linear-Riemann-piece over $(\lambda,\mu,\nu)$ vanishes:
\begin{equation}
\partial_{\lambda} \big(\partial_{\mu} \Christoffel^{\rho}{}_{\nu\sigma}
                       - \partial_{\nu} \Christoffel^{\rho}{}_{\mu\sigma}\big)
+ \partial_{\mu} \big(\partial_{\nu} \Christoffel^{\rho}{}_{\lambda\sigma}
                       - \partial_{\lambda} \Christoffel^{\rho}{}_{\nu\sigma}\big)
+ \partial_{\nu} \big(\partial_{\lambda} \Christoffel^{\rho}{}_{\mu\sigma}
                       - \partial_{\mu} \Christoffel^{\rho}{}_{\lambda\sigma}\big)
= 0.
\label{eq:linear-diff-bianchi}
\end{equation}
\end{theorem}

This is the substantive Schwarz-only discharge. Three pairs of
$d^{2}\Christoffel$ summands cancel exactly under
$d^{2}\Christoffel_{\lambda\mu} = d^{2}\Christoffel_{\mu\lambda}$,
$d^{2}\Christoffel_{\mu\nu} = d^{2}\Christoffel_{\nu\mu}$, and
$d^{2}\Christoffel_{\nu\lambda} = d^{2}\Christoffel_{\lambda\nu}$. No
torsion-freeness on $\Christoffel$ or first Bianchi on $R$ is
needed for the linear piece — it is the only sub-piece of the
differential Bianchi that decouples cleanly at coordinate scope.

\subsection{Group C+D quadratic cancellation}
\label{sec:groupcd}
The $\Christoffel \times R$ correction terms in
\eqref{eq:cov-riemann} split into four groups by which index of $R$
the contraction-with-$\Christoffel$ index $\alpha$ lands on. Under
the cyclic sum in $(\lambda, \mu, \nu)$, the lower-index pair of
groups (the $\mu$- and $\nu$-position contractions, denoted Groups C
and D) cancels by torsion-freeness on $\Christoffel$ plus
\texttt{AntisymLastTwo} on $R$ alone:

\begin{theorem}[Group C+D cyclic cancellation]
\label{thm:group-cd}
For any $\Christoffel$ with \texttt{IsTorsionFree} and any $R$ with
\texttt{AntisymLastTwo},
\begin{equation}
\sum_{\mathrm{cyclic}\,(\lambda,\mu,\nu)}
  \Big(- \sum_{\alpha} \Christoffel^{\alpha}{}_{\lambda\mu}
        R^{\rho}{}_{\sigma\alpha\nu}
       - \sum_{\alpha} \Christoffel^{\alpha}{}_{\lambda\nu}
        R^{\rho}{}_{\sigma\mu\alpha}\Big) = 0.
\label{eq:group-cd}
\end{equation}
\end{theorem}

The proof is a 48-monomial signed-pair cancellation: per-$\alpha$,
$\Christoffel^{\alpha}{}_{\lambda\mu} = \Christoffel^{\alpha}{}_{\mu\lambda}$
(torsion-free) and $R^{\rho}{}_{\sigma\alpha\nu}
= -R^{\rho}{}_{\sigma\nu\alpha}$ (antisym), giving the
inter-cyclic match across all four index values $\alpha \in \{0,1,2,3\}$
in each of the six summands. This is the cleanest non-trivial chunk of
the full differential Bianchi that holds at coordinate scope.

\subsection{Strategic deferral of the full discharge to Session 4}
\label{sec:deferral}
The cyclic sum of $\nabla_{\lambda} R^{\rho}{}_{\sigma\mu\nu}$ over
$(\lambda, \mu, \nu)$ decomposes into four pieces, each requiring
distinct hypothesis combinations:

\begin{itemize}
\item $(\partial^{2}\Christoffel)$-piece: cancels under \texttt{IsSchwarz}
  alone (Theorem \ref{thm:linear-diff-bianchi}).
\item $(\partial \Christoffel \times \Christoffel)$ bilinear piece:
  needs cross-pair matching between $\partial(\Christoffel\Christoffel)$
  and Group A/B $(\Christoffel \times \partial\Christoffel)$ terms;
  cancellation requires algebraic first Bianchi on $R$ + torsion-freeness
  on $\Christoffel$.
\item $(\Christoffel\Christoffel\Christoffel)$ cubic piece: cancels
  under torsion-free + algebraic first Bianchi via standard
  symmetric-pair matching.
\item Group C+D lower-index correction: cancels under torsion-free +
  \texttt{AntisymLastTwo} alone (Theorem \ref{thm:group-cd}).
\end{itemize}

The bilinear and cubic parts are computationally tractable at coordinate
scope but produce hundreds of \texttt{ring}-normalised monomials. At
bundle scope (Wave 8 Session 3, the Bonn
\texttt{IsCovariantDerivativeOn} consumer; see \S\ref{sec:bundle-aux}
below for the Session 3 module that opens that scope), the same content
discharges in three lines via $\texttt{mlieBracket\_jacobi}$. The
strategic ordering of the five-session plan trades a one-session delay
for a ten-fold simpler proof, and the Session 4 specialisation theorem
then delivers the full coordinate-level identity as a consequence.

\section{Bundle-level Riemann curvature on the Bonn API}\label{sec:bundle-aux}

The bridge from coordinate-scope (Sessions 1+2) to bundle-scope
(Sessions 3--5) is opened by the Session 3 module
\texttt{BundleRiemannAux.lean}, which lifts the curvature operator
from the algebraic-precedent $\mathrm{Fin}\,4 \to \mathrm{Fin}\,4 \to
\mathrm{Fin}\,4 \to \mathrm{Fin}\,4 \to \mathbb{R}$ data array to a
function $(\mathrm{cov}, X, Y, Z, x) \mapsto V\,x$ on a vector bundle
$V \to M$ over a manifold $M$ over a nontrivially-normed field
$\mathbb{K}$. The curvature operator is defined verbatim from
Klingenberg's textbook formulation \cite{Klingenberg1995},
\begin{equation}
R(X, Y) Z := \nabla_{X} \nabla_{Y} Z - \nabla_{Y} \nabla_{X} Z
            - \nabla_{[X, Y]} Z,
\label{eq:bundle-riem}
\end{equation}
unrolled into the Bonn-API argument-order convention
$\mathrm{cov}\,\sigma\,x\,(X\,x) \equiv \nabla_{X}\,\sigma\,x$:
\begin{align}
\mathtt{bundleRiemannAux}\,\mathrm{cov}\,X\,Y\,Z\,x
\;=\;
&\,\mathrm{cov}(\,y \mapsto \mathrm{cov}\,Z\,y\,(Y\,y)\,)\,x\,(X\,x)
\nonumber \\
& \!\! - \mathrm{cov}(\,y \mapsto \mathrm{cov}\,Z\,y\,(X\,y)\,)\,x\,(Y\,x)
\nonumber \\
& \!\! - \mathrm{cov}\,Z\,x\,(\mathrm{mlieBracket}\,I\,X\,Y\,x).
\label{eq:bundle-riem-aux}
\end{align}

\subsection{Wave-headlines: $R(X, X)Z = 0$ and antisymmetry-$(X, Y)$}
\label{sec:bundle-aux-headlines}
The two wave-headlines of Session 3 are purely algebraic consequences
of (a) the manifold Lie bracket lemmas
$\texttt{VectorField.mlieBracket\_self}$ and
$\texttt{VectorField.mlieBracket\_swap\_apply}$ in
\cite{Gouezel2024LieBracket}, and (b) the type-level
continuous-linearity of $\mathrm{cov}\,\sigma\,x : T_xM \to_L V_x$
($\mathtt{map\_zero}$ and $\mathtt{map\_neg}$). They do \emph{not}
require $\mathrm{cov}$ to be a covariant derivative.

\begin{theorem}[$R(X, X)Z = 0$]\label{thm:bundleRiem-self}
For any $\mathrm{cov} : (\Pi\,x : M, V\,x) \to (\Pi\,x : M, T_xM \to_L V_x)$,
any vector field $X$, any section $Z$, and any point $x$,
$\mathtt{bundleRiemannAux}\,\mathrm{cov}\,X\,X\,Z\,x = 0$.
\end{theorem}

The proof unfolds the definition; the first two summands cancel
immediately ($A - A = 0$); the third summand vanishes via
$\mathtt{mlieBracket\_self} : \mathrm{mlieBracket}\,I\,X\,X = 0$
(at point $x$, $0 : T_xM$) followed by
$(\mathrm{cov}\,Z\,x).\mathtt{map\_zero}$.

\begin{theorem}[Antisymmetry-$(X, Y)$ at a point]\label{thm:bundleRiem-swap}
For any $\mathrm{cov}$, any vector fields $X, Y$, any section $Z$,
and any point $x$,
$\mathtt{bundleRiemannAux}\,\mathrm{cov}\,X\,Y\,Z\,x
  = -\,\mathtt{bundleRiemannAux}\,\mathrm{cov}\,Y\,X\,Z\,x$.
\end{theorem}

The proof rewrites the bracket via
$\mathtt{mlieBracket\_swap\_apply} : \mathrm{mlieBracket}\,I\,X\,Y\,x
  = -\,\mathrm{mlieBracket}\,I\,Y\,X\,x$, then linearity of
$\mathrm{cov}\,Z\,x$ converts $\mathrm{cov}\,Z\,x\,(-v) =
-\,\mathrm{cov}\,Z\,x\,v$, and the four leading summands swap and
negate to match. The function-equality companion
$\texttt{bundleRiemannAux\_swap}$ is shipped via $\mathtt{funext}$ in
the Mathlib idiom of pairing $\texttt{\_apply}$ with its section-level
form.

\subsection{Commuting-vector-fields specialisation}
When the Lie bracket vanishes at a point (e.g.\ when $X$ and $Y$ are
coordinate vector fields, or commute by some structural reason), the
curvature reduces to the pure connection commutator. We ship this as
a substantive named simplification:

\begin{theorem}[Commuting-vector-fields specialisation]\label{thm:bundleRiem-commute}
Under the hypothesis $\mathrm{mlieBracket}\,I\,X\,Y\,x = 0$,
$\mathtt{bundleRiemannAux}\,\mathrm{cov}\,X\,Y\,Z\,x =
\mathrm{cov}(\,y \mapsto \mathrm{cov}\,Z\,y\,(Y\,y)\,)\,x\,(X\,x) -
\mathrm{cov}(\,y \mapsto \mathrm{cov}\,Z\,y\,(X\,y)\,)\,x\,(Y\,x)$.
\end{theorem}

This is the form encountered in textbook coordinate calculations
(where $X = \partial_{\mu}, Y = \partial_{\nu}$ have vanishing
bracket) and is the structural cross-bridge to the Sessions-1+2
coordinate Riemann that Session 4's specialisation theorem will
upgrade.

\subsection{Zero-degeneracies on Bonn's
\texttt{IsCovariantDerivativeOn}}
\label{sec:bundle-aux-zero}
Two further structural identities ship under the hypothesis
$\mathtt{IsCovariantDerivativeOn}\,F\,\mathrm{cov}\,\mathrm{Set.univ}$
and the typeclass $[\mathtt{VectorBundle}\,\mathbb{K}\,F\,V]$, which
together supply $\mathrm{cov}\,0 = 0$ via Bonn's
$\texttt{IsCovariantDerivativeOn.zero}$ \cite{MassotRothgangMacbeth2025}.

\begin{theorem}[$R(X, Y)\,0 = 0$]\label{thm:bundleRiem-zero-section}
$\mathtt{bundleRiemannAux}\,\mathrm{cov}\,X\,Y\,(0 : \Pi\,x, V\,x)\,x
  = 0$, with the proof composing $\mathrm{cov}\,0 = 0$ at the inner
sections $(\,y \mapsto \mathrm{cov}\,0\,y\,(Y\,y)\,) = 0$ and
$(\,y \mapsto \mathrm{cov}\,0\,y\,(X\,y)\,) = 0$, and at the outer
$\mathrm{cov}\,0\,x = 0$.
\end{theorem}

\begin{theorem}[$R(0, Y)\,Z = 0$]\label{thm:bundleRiem-zero-X}
$\mathtt{bundleRiemannAux}\,\mathrm{cov}\,(0 : \Pi\,x, T_xM)\,Y\,Z\,x
  = 0$, composing four distinct Mathlib lemmas:
$\mathtt{Pi.zero\_apply}$ for the outer X-slot;
$(\mathrm{cov}\,\sigma\,x).\mathtt{map\_zero}$;
$(\,y \mapsto \mathrm{cov}\,Z\,y\,0\,) = 0$ followed by
$\mathtt{IsCovariantDerivativeOn.zero}$;
$\mathtt{mlieBracket\_zero\_left}$ for the bracket term.
\end{theorem}

The dual $R(X, 0)\,Z = 0$ falls out as a one-line corollary of
Theorem~\ref{thm:bundleRiem-swap} composed with
Theorem~\ref{thm:bundleRiem-zero-X} and is intentionally not shipped
to avoid P3-trivial corollary plumbing.

\subsection{Why bundle scope (and what is intentionally deferred to
Session 4)}\label{sec:bundle-aux-defer}
The Bianchi identities at bundle scope follow with much less arithmetic
than at coordinate scope because they reduce to algebraic properties
of the Lie bracket. The algebraic first Bianchi at bundle scope for
torsion-free $\mathrm{cov}$, $R(X, Y) Z + R(Y, Z) X + R(Z, X) Y = 0$,
follows from torsion vanishing plus the Jacobi identity for
$\mathrm{mlieBracket}$ in three lines. The differential second Bianchi
at bundle scope, $(\nabla_X R)(Y, Z) + (\nabla_Y R)(Z, X) + (\nabla_Z R)
(X, Y) = 0$, falls out from Jacobi for $\mathrm{mlieBracket}$ and from
$(\nabla_X R)$ being itself a tensor, also in three lines.

These payoffs make Session 4's specialisation-from-bundle the right
architecture: at coordinate scope the same identities require $\sim 500$
monomials of \texttt{ring}-discharge after dozens of explicit rewrites
(per Sessions 1+2 strategic deferrals). Session 4 will (a)~build
$\texttt{TensorialAt.mkHom}_3$ (Bonn already has $\mathtt{mkHom}_2$
for \texttt{torsion}), (b)~bundle $\mathtt{bundleRiemannAux}$ into a
$\texttt{CovariantDerivative.riemann}$ whose value at each fiber is a
trilinear $V_x \to_L V_x \to_L V_x \to_L V_x$ map, (c)~ship the
bundle-level first and second Bianchi as the three-line Jacobi
consequences, and (d)~specialise on $\mathrm{EuclideanSpace}\,
\mathbb{R}\,(\mathrm{Fin}\,4)$ to deliver the Sessions-1+2
coordinate-level identities together with the load-bearing
$\nabla^{\mu}G_{\mu\nu} = 0$ discharge for
\texttt{EinsteinTensor.lean}.

\section{What this unblocks}\label{sec:unblocks}

The four new modules are infrastructure: they ship no new physics and
no new submission-ready prediction. Their value is in unblocking five
deferred items in the Phase 6f programme that were previously parked
behind the Mathlib indefinite-signature gap and the
Riemann-from-connection gap. The W8 Session 1 first-Bianchi discharge
in particular promotes the unblock from ``mechanical pending Mathlib"
to ``mechanical pending an internal SK-EFT module" for several items:

\begin{itemize}
\item \textbf{6f.2 second Bianchi.}
  $\nabla^{\mu}G_{\mu\nu} = 0$ requires a covariant-divergence
  operator on the Einstein tensor. With \texttt{christoffelKoszul} in
  hand and the algebraic first Bianchi shipped, the divergence can
  be expressed as an algebraic identity in
  $(\MetricMatrix,\PartG,\partial^{2}\MetricMatrix)$ and the second
  Bianchi proven directly. The W8 Session 2 module
  \texttt{RiemannDifferentialBianchi.lean} ships the
  $\partial^{2}\Christoffel$ carrier with the Schwarz hypothesis, the
  covariant-derivative operator on a generic Riemann tensor, the
  linear-piece differential second Bianchi under Schwarz alone, and
  the Group C+D quadratic cyclic cancellation under torsion-freeness
  plus \texttt{AntisymLastTwo}. The full coordinate-level differential
  second Bianchi for the entire Riemann tensor is deferred to the W8
  Session 4 specialisation theorem, which delivers it as a consequence
  of the Session 3 bundle-level Jacobi identity. The remaining gap
  between the formalisation and \texttt{EinsteinTensor.lean}'s
  $\nabla^{\mu}G_{\mu\nu}=0$ theorem statement is therefore the
  Session 4 specialisation; both layers (bilinear and cubic) of the
  full discharge are themselves three-line bundle-level statements.
\item \textbf{6f.4 Schwarzschild vacuum-Ricci.}
  The algebraic-precedent statement $R_{\mu\nu} = 0$ for the
  Schwarzschild solution was rejected at first pass per the
  6f.2 lesson (the tracked-Prop form was vacuous because we had no
  Christoffel-side machinery). With the full coordinate Riemann
  including the linear-in-$\partial\Christoffel$ piece shipped, the
  $\Riem^{\rho}{}_{\sigma\rho\nu} = 0$ contraction at the algebraic
  level is now mechanical: instantiate $\Christoffel$ and
  $\partial\Christoffel$ to the Schwarzschild values, contract over
  the upper Riemann index, and discharge by direct computation. This
  is naturally a follow-on consumer module after Session 2.
\item \textbf{6f.5 ADM 3D Riemannian Ricci on Cauchy slices.}
  The $D_{j}$ covariant-divergence on a 3D Cauchy slice was deferred
  for the same reason. With the full coordinate Riemann shipped,
  the 3D restriction is direct: instantiate \texttt{Christoffel} to
  the spatial connection and reuse the same algebraic-precedent
  pattern.
\item \textbf{6f.6 Cartan structure equations.}
  $d\omega + \omega \wedge \omega = R$ requires the connection 1-form
  $\omega$, which is naturally read off the Christoffel symbols via
  $\omega^{a}{}_{b\mu} = e^{a}{}_{\nu}\Christoffel^{\nu}{}_{b\mu}
  - e^{a}{}_{b,\mu}$. With Christoffel symbols formalized and the
  algebraic first Bianchi corollary shipped (which gives the
  algebraic identity $d^2 \omega = 0$ in the cyclic-sum form), the
  bridge is a tetrad/connection-1-form arithmetic exercise once the
  exterior derivative $d$ is modelled (W8 Session 4 bundle-level
  layer is the natural place; or a separate
  \texttt{DifferentialFormsCoordinate} sibling module).
\item \textbf{6g.2 / 6g.6 curve-theoretic Penrose / Mazur.}
  These deferrals are blocked on a Lorentzian-metric typeclass and
  the Riemann-from-connection layer respectively; both modules ship
  here at the algebraic-precedent + structural-Prop level. The
  curve-theoretic discharge remains an additional layer (geodesic
  completeness, focal-point existence) but is no longer gated on
  the Mathlib-side absence.
\end{itemize}

\section{Conventions and conformance}\label{sec:conventions}

The five new Lean modules follow Mathlib upstream conventions: file
header copyright + module imports + \texttt{public} sections, no
\texttt{maxHeartbeats} overrides, no \texttt{ring} on noncommutative
contexts. Naming follows Mathlib snake-case for theorems and
PascalCase for predicates and structures. The Session 3 module
\texttt{BundleRiemannAux.lean} additionally adopts Bonn's variable
section verbatim (model-with-corners $I$, charted-space $M$, vector
bundle $V \to M$ with model fiber $F$), uses \texttt{omit} clauses on
the algebraic theorems that do not consume
$[\forall x, \mathrm{IsTopologicalAddGroup}\,(V\,x)]$ /
$[\forall x, \mathrm{ContinuousSMul}\,\mathbb{K}\,(V\,x)]$ in the
Mathlib upstream-style idiom (cf.\ \texttt{Mathlib.Geometry.Manifold.\-
ContMDiff.Atlas} for parallel \texttt{omit} usage), and opens
\texttt{VectorField} so that \texttt{mlieBracket\_self} and
\texttt{mlieBracket\_swap\_apply} resolve directly without
fully-qualified names. Cross-module bridges from
the \texttt{LorentzianMetric} module reach into existing
\texttt{LinearizedEFE.$\eta$\_symm}; the
\texttt{RiemannianConnection} module reaches into
\texttt{Curvature.AntisymLastTwo}; the
\texttt{RiemannCoordinate} module reaches into both
\texttt{Curvature.FirstBianchi} (the wave-headline first-Bianchi
predicate consumer) and W7's
\texttt{RiemannianConnection.IsTorsionFree} +
\texttt{RiemannianConnection.riemannQuadraticTerm}. All cross-module
calls are SK-EFT-internal and survive the eventual upstream port —
they will be replaced by their Mathlib-bundle-level counterparts when
the \texttt{IsContinuousLorentzianBundle} typeclass and the
\texttt{CovariantDerivative}-consuming Riemann tensor land upstream.

The library passes the project's preemptive-strengthening discipline
\cite{Roehm2026Pipeline}: each theorem statement is non-vacuous under
its hypothesis, no $\exists$-absorption, no bundle-redundancy, and
the Riemannian-Lorentzian falsifier and the
metric-partial-symmetry-load-bearing Koszul torsion-freeness
preempt P3/P5-trivial restatements.

\subsection{Strengthening cuts}\label{sec:strengthening}

\paragraph{W7 first-pass cut.}
The W7 wave applied a single retroactive cut at first-pass review:
\texttt{zeroChristoffel\_isTorsionFree}, a pure-\texttt{rfl} P3-trivial
witness with no downstream consumer (zero-symmetry follows by
definition without consuming any hypothesis).

\paragraph{W8 Session-1+2+3 first-pass cuts.}
W8 Sessions 1, 2, and 3 each shipped with zero retroactive cuts at
first pass and after the ruthless audit — matching the project
best-pass discipline metric (6e.5 / 6b.2 / 6f.6 baselines).

\paragraph{W8 post-Session-3 strengthening pass (two retroactive cuts).}
A reflexive end-of-session strengthening pass after Session 3 found
two cut candidates inherited from W7 that the W7 first-pass review
missed:

\begin{enumerate}
  \item \texttt{IsLeviCivitaForKoszul} predicate def +
    \texttt{leviCivita\_unique\_at\_koszul} theorem, both in
    \texttt{RiemannianConnection.lean}. The predicate
    \texttt{IsLeviCivitaForKoszul $\Christoffel\, \MetricMatrix^{-1}\,
    \PartG := \Christoffel = \texttt{christoffelKoszul}\,
    \MetricMatrix^{-1}\,\PartG$} is literally the equality
    "$\Christoffel$ equals the closed-form Koszul Christoffel". The
    "uniqueness" theorem proof was \texttt{rw [h, h']} — pure
    transitivity of equality through definitional unfolding. This is
    a P5 trivial-trans on a def-as-equality predicate. The substantive
    Levi-Civita uniqueness — the Wald §3.1 / Kobayashi-Nomizu Thm
    IV.2.2 fundamental theorem that every torsion-free
    metric-compatible connection equals the Koszul Christoffel — needs
    the bundle-level Koszul-bilinear-form argument over arbitrary
    smooth vector fields and is the Wave 8 Session 5 deliverable.
    Documented in Section~\ref{sec:koszul} Remark~\ref{rem:lc-unique-deferred}.
  \item \texttt{lorentzian\_metric\_nonzero} theorem in
    \texttt{LorentzianMetric.lean}. The conclusion
    "$\MetricMatrix\,0\,0 \neq 0$" follows from the
    \texttt{IsLorentzianSignature} predicate's first conjunct
    "$\MetricMatrix\,0\,0 = -1$" by one-step \texttt{rw + norm\_num}
    discharge of $-1 \neq 0$. The substantive non-vacuity content
    lives in the \texttt{IsLorentzianSignature} predicate definition
    itself; the determinant-discriminator content lives in
    \texttt{lorentzian\_diagonal\_product\_neg\_one}. No downstream
    consumer.
\end{enumerate}

Both cuts retired predicate-level trivial-trans / corollary patterns
that had no downstream consumers and whose substantive content lives
elsewhere in the library. After the cuts the library ships
\textbf{thirty-three substantive theorems across five modules},
zero \texttt{sorry}, zero new project axioms; all load-bearing
theorems verify against \texttt{propext, Classical.choice, Quot.sound}
only. The pattern lesson is logged in
\texttt{feedback\_post\_wave\_strengthening\_audit.md}: predicate-as-
equality definitions \emph{should not} appear with companion
"uniqueness" theorems at the algebraic-precedent scope; defer
substantive uniqueness to the bundle-level wave where the Koszul
bilinear-form argument supplies the load-bearing content.

\section{Discussion and follow-up}\label{sec:disc}

\subsection{Bundle-level Mathlib upstream port}
The natural upstream-port target is a \texttt{Curvature.lean} file
under \texttt{Mathlib.Geometry.Manifold.VectorBundle.CovariantDerivative}
that defines the Riemann curvature on a bundled \texttt{CovariantDerivative}
via the connection-commutator function
$\Riem(X,Y)Z := \nabla_{X}(\nabla_{Y}Z) - \nabla_{Y}(\nabla_{X}Z)
- \nabla_{[X,Y]}Z$, proves antisymmetry-in-$(X,Y)$ via the
\texttt{TensorialAt.mkHom\textsubscript{2}} machinery already used in
\texttt{IsCovariantDerivativeOn.torsion}, then derives Ricci, scalar
curvature, the algebraic Bianchi identity (under torsion-free
hypothesis), and the bundle-level Levi-Civita uniqueness via the
Koszul-bilinear-form argument over smooth vector fields. The
algebraic-precedent identities formalized here are the substantive
inputs that survive the bundle-level upgrade.

A parallel \texttt{Lorentzian.lean} sibling of
\texttt{Mathlib.Topology.VectorBundle.Riemannian} would supply the
\texttt{IsContinuousLorentzianBundle} typeclass shape; the
algebraic-precedent witness here lifts directly.

\subsection{Comparison with HoTT / HOL Light / Coq prior art}
We are unaware of a Levi-Civita Christoffel formalization in any
proof assistant. HOL Light's geometric-algebra library provides
Riemannian inner-product spaces but no Christoffel/connection.
Isabelle's AFP has nothing in this space. Coq's MathComp has
neither Riemannian nor Lorentzian metric infrastructure. PhysLean
ports physics formulas but has no curvature library. The HALF-ERC
audit report's no-prior-art finding therefore stands.

\subsection{Multi-session plan and open work}
The bundle-level upgrade is now an explicit five-session plan tracked
in the Phase 6f roadmap (Wave 8, sessions 1--5). Sessions 1, 2, and 3
ship the content reported in this paper.
Session 2 adds \texttt{RiemannDifferentialBianchi.lean} with
$\partial^{2}\Christoffel$ + Schwarz symmetry + the linear-piece
differential second Bianchi identity
$\nabla_{\lambda}\Riem + \nabla_{\mu}\Riem + \nabla_{\nu}\Riem = 0$,
which mechanically discharges the
\texttt{EinsteinTensor.}$\nabla^{\mu}G_{\mu\nu} = 0$ deferral.
Session 3 (\S\ref{sec:bundle-aux}) adds \texttt{BundleRiemannAux.lean}
opening the Bonn-API bundle scope: the bare-function
\texttt{bundleRiemannAux} consuming Bonn
\texttt{IsCovariantDerivativeOn} and \texttt{VectorField.mlieBracket}
(antisymmetric in $(X,Y)$ from $\mathrm{mlieBracket}$ antisymmetry,
vanishing on the diagonal from $\mathtt{mlieBracket\_self}$, with the
commuting-vector-fields specialisation and the section/X-slot
zero-degeneracies on $\texttt{IsCovariantDerivativeOn.zero}$ shipped
in the same module). Session 4 builds tensoriality discharges via
\texttt{TensorialAt.mkHom\textsubscript{3}} (which we will build from
\texttt{mkHom\textsubscript{2}} if not already in Mathlib), the
bundled \texttt{CovariantDerivative.riemann}, the bundle-level
algebraic first Bianchi for torsion-free (a three-line consequence of
the Jacobi identity for $\mathrm{mlieBracket}$), the bundle-level
differential second Bianchi (also three-line via Jacobi), and a
specialisation theorem bridging the bundle Riemann on
\texttt{EuclideanSpace}\,$\mathbb{R}\,(\mathrm{Fin}\,4)$ to the
coordinate Riemann from Sessions 1 and 2 (which delivers the full
coordinate-level differential second Bianchi at top order plus the
$\nabla^{\mu}G_{\mu\nu} = 0$ discharge as direct consequences).
Session 5 ships Levi-Civita uniqueness on a vector bundle (Koszul
bilinear form on smooth sections), sends the staged Bonn-team
outreach email per the project's
\texttt{phase6f1\_curvature\_architecture\_decision.md} script, and
submits the Mathlib pull request. Each session updates the running
state checklist in the Phase 6f roadmap and ships a session-specific
project memory; the auto-compact-survival contract is documented in
the roadmap's Wave 8 ``Session-handoff" subsection.

The full curve-theoretic Penrose and Mazur-Bunting-Israel proofs,
also blocked, have additional gates beyond the Lorentzian-metric
availability (geodesic existence, focal-point existence,
$\sigma$-model rigidity); those are phase-level commitments. Within
the SK-EFT-Hawking programme they sit alongside the Phase 6g.7
cosmic-censorship backlog.

\section*{Acknowledgments}
We thank the Mathlib community, particularly Patrick Massot, Michael
Rothgang, and Heather Macbeth for the covariant-derivative library;
S\'ebastien Gou\"ezel for the Riemannian-bundle infrastructure that
the Lorentzian companion mirrors; and the SK-EFT-Hawking project's
preemptive-strengthening discipline for catching the
$\mathrm{zeroChristoffel\_isTorsionFree}$ P3-trivial cut at the
post-wave review.

\bibliographystyle{apsrev4-2}
\begin{thebibliography}{99}

\bibitem{MassotRothgangMacbeth2025}
P.~Massot, M.~Rothgang, H.~Macbeth.
\newblock Mathlib4 covariant derivatives library.
\newblock \emph{Mathlib4 commit 8850ed93}, April 2026.
\newblock \texttt{Mathlib.Geometry.Manifold.VectorBundle.CovariantDerivative.\{Basic,Torsion\}}.

\bibitem{Gouezel2025}
S.~Gou\"ezel.
\newblock Riemannian vector bundles in Mathlib4.
\newblock \emph{Mathlib4 commit 8850ed93}, March 2026.
\newblock \texttt{Mathlib.Topology.VectorBundle.Riemannian};
  \texttt{Mathlib.Geometry.Manifold.VectorBundle.Riemannian}.

\bibitem{Roehm2026Wave6f}
J.\,G.~Roehm.
\newblock Phase 6f waves 1--6 module roster: classical-GR algebraic
  backbone in Lean.
\newblock \emph{SK-EFT-Hawking project, in preparation},
  inventory snapshot 2026-04-29.

\bibitem{Roehm2026Wave6g}
J.\,G.~Roehm.
\newblock Phase 6g waves 1--6: causal structure, singularity
  hypothesis bundles, area theorem, no-hair theorem in Lean.
\newblock \emph{SK-EFT-Hawking project, in preparation},
  inventory snapshot 2026-04-29.

\bibitem{Roehm2026Pipeline}
J.\,G.~Roehm.
\newblock Wave-execution pipeline and preemptive-strengthening
  discipline for theorem-driven theoretical-physics workflows.
\newblock \emph{SK-EFT-Hawking project, in preparation},
  \texttt{docs/WAVE\_EXECUTION\_PIPELINE.md}.

\bibitem{Wald1984}
R.\,M.~Wald.
\newblock \emph{General Relativity}.
\newblock University of Chicago Press, 1984.
\newblock \S 3.1, 3.2.

\bibitem{KobayashiNomizu1963}
S.~Kobayashi and K.~Nomizu.
\newblock \emph{Foundations of Differential Geometry, Vol.~I}.
\newblock Wiley-Interscience, 1963.
\newblock Theorem IV.2.2 (Levi-Civita uniqueness).

\bibitem{Carroll2004}
S.~Carroll.
\newblock \emph{Spacetime and Geometry}.
\newblock Addison-Wesley, 2004.
\newblock \S 3.2 (Koszul formula).

\bibitem{Klingenberg1995}
W.~Klingenberg.
\newblock \emph{Riemannian Geometry}, 2nd ed.
\newblock de Gruyter Studies in Mathematics, vol.\ 1, 1995.
\newblock \S 1.3 (curvature operator on a vector bundle).

\bibitem{Gouezel2024LieBracket}
S.~Gou\"ezel.
\newblock Lie brackets of vector fields on manifolds in Mathlib4.
\newblock \emph{Mathlib4 commit 8850ed93}, 2024--2026.
\newblock \texttt{Mathlib.Geometry.Manifold.VectorField.LieBracket}
  (\texttt{mlieBracket}, \texttt{mlieBracket\_self},
  \texttt{mlieBracket\_swap\_apply}, \texttt{mlieBracket\_zero\_left},
  \texttt{leibniz\_identity\_mlieBracket}).

\end{thebibliography}

\end{document}
